\section{Bölüm sonu çözümlü problemler}

\begin{enumerate}
  \item \textbf{Denklem ve tahmin.} $\widehat Y=70-2{,}5X$ modelinde $X=8$ için tahmini bulun ve eğimi yorumlayın.

  \textbf{Çözüm:} $\widehat Y=70-2{,}5(8)=50$'dir. X'teki bir birimlik artış, beklenen Y ortalamasında 2,5 birimlik azalmayla ilişkilidir. Gözlemsel veride bu ifade müdahale etkisi değildir.

  \item \textbf{Artık ve uyum.} Bir öğrencinin gözlenen puanı 58, model tahmini 52 ise artığı bulun. $SST=900$ ve $SSE=540$ için $R^2$ değerini hesaplayın.

  \textbf{Çözüm:} Artık $e=58-52=6$'dır; model öğrencinin puanını 6 birim düşük tahmin etmiştir. $R^2=1-SSE/SST=1-540/900=0{,}40$ olur.

  \item \textbf{İki aralığı ayırma.} $X=6$ için ortalama yanıt güven aralığı $[45;51]$, bireysel tahmin aralığı $[32;64]$ bulunmuştur. Farkı açıklayın.

  \textbf{Çözüm:} İlk aralık, $X=6$ olan birimlerin evren ortalamasını hedefler. İkinci aralık yeni tek bir bireyin doğal artık değişkenliğini de içerir; bu nedenle daha geniştir. Tek öğrencinin puanı ortalama güven aralığıyla değerlendirilmemelidir.
\end{enumerate}